\documentclass[11pt]{article}

\usepackage[T1]{fontenc}
\usepackage[utf8]{inputenc}
\usepackage[a4paper,margin=1in]{geometry}
\usepackage{amsmath,amssymb}
\usepackage{booktabs}
\usepackage{hyperref}
\usepackage{xcolor}

\hypersetup{
  colorlinks=true,
  linkcolor=blue,
  citecolor=blue,
  urlcolor=blue
}

\title{Pilot Methodology for Meyer--Wallach Entanglement Evaluation on IQM Spark}
\author{QC1 Quantum Banknote Classifier on ODRA 5}
\date{}

\begin{document}

\maketitle

\section{Purpose}

This document specifies a pilot protocol for choosing the hardware budget in the
Meyer--Wallach (MW) entanglement workflow. The protocol selects:
\begin{itemize}
  \item the number of shots per Pauli measurement basis,
  \item the number of random ansatz parameter samples,
  \item and the number of repeated fixed-seed hardware iterations used to
  quantify run-to-run drift.
\end{itemize}

The corresponding runner is:
\begin{verbatim}
python scripts/run_iqm_mw_pilot.py --pilot-id mw_pilot_paper
\end{verbatim}

The output directory is:
\begin{verbatim}
evaluation_and_comparison/iqm_spark/iqm_mw_outputs/pilots/<pilot_id>/
\end{verbatim}

\section{Quantity Being Estimated}

For an \(n\)-qubit pure state, the Meyer--Wallach score is
\begin{equation}
Q(|\psi\rangle)
=
\frac{2}{n}
\sum_{i=1}^{n}
\left(1-\operatorname{Tr}(\rho_i^2)\right),
\label{eq:mw}
\end{equation}
where \(\rho_i\) is the one-qubit reduced density matrix of qubit \(i\).

The hardware workflow does not reconstruct the full \(2^n \times 2^n\) density
matrix. Instead, it estimates only the one-qubit marginals required by
Eq.~\eqref{eq:mw}. Each reduced state is represented by its Bloch vector,
\begin{equation}
\rho_i
=
\frac{1}{2}
\left(I+x_iX+y_iY+z_iZ\right),
\label{eq:bloch}
\end{equation}
where
\begin{equation}
x_i=\langle X_i\rangle,\qquad
y_i=\langle Y_i\rangle,\qquad
z_i=\langle Z_i\rangle.
\end{equation}

The single-qubit purity is therefore
\begin{equation}
\operatorname{Tr}(\rho_i^2)
=
\frac{1}{2}
\left(1+x_i^2+y_i^2+z_i^2\right).
\label{eq:purity}
\end{equation}
Substituting Eq.~\eqref{eq:purity} into Eq.~\eqref{eq:mw} gives
\begin{equation}
Q
=
1-\frac{1}{n}
\sum_{i=1}^{n}
\left(x_i^2+y_i^2+z_i^2\right).
\label{eq:mw_bloch}
\end{equation}

Thus, for each random ansatz parameter sample, the hardware run needs only
three circuits: one measured in the \(Z\) basis, one in the \(X\) basis, and one
in the \(Y\) basis.

\section{Shot-Count Justification}

For a fixed qubit \(i\) and Pauli basis \(P \in \{X,Y,Z\}\), each shot returns
a binary outcome
\begin{equation}
Y_s \in \{-1,+1\}.
\end{equation}
The Pauli expectation estimator is the sample mean
\begin{equation}
\hat p_{iP}
=
\frac{1}{S}
\sum_{s=1}^{S} Y_s,
\end{equation}
where \(S\) is the number of shots in that basis. If
\(p_{iP}=\mathbb{E}[Y_s]\), then \(Y_s^2=1\), so
\begin{equation}
\operatorname{Var}(Y_s)
=
\mathbb{E}[Y_s^2]-\mathbb{E}[Y_s]^2
=
1-p_{iP}^2.
\end{equation}
For independent shots,
\begin{equation}
\operatorname{Var}(\hat p_{iP})
=
\frac{1-p_{iP}^2}{S}
\leq
\frac{1}{S}.
\label{eq:pauli_var}
\end{equation}

Since
\begin{equation}
Q
=
1-\frac{1}{n}\sum_{i,P}p_{iP}^2,
\end{equation}
the delta method gives the approximate propagation of Pauli-estimation shot
noise into one MW score:
\begin{equation}
\operatorname{Var}(\hat Q)
\approx
\sum_{i,P}
\left(
\frac{\partial Q}{\partial p_{iP}}
\right)^2
\operatorname{Var}(\hat p_{iP})
=
\sum_{i,P}
\left(
\frac{2p_{iP}}{n}
\right)^2
\frac{1-p_{iP}^2}{S}.
\label{eq:delta}
\end{equation}

Using the Bloch-vector constraint
\begin{equation}
x_i^2+y_i^2+z_i^2 \leq 1,
\end{equation}
Eq.~\eqref{eq:delta} is conservatively bounded by
\begin{equation}
\operatorname{SD}_{\mathrm{shot}}(\hat Q)
\lesssim
\frac{2}{\sqrt{nS}}.
\label{eq:shot_bound}
\end{equation}

For the five-qubit IQM Spark experiment, \(n=5\). With \(S=4096\) shots,
\begin{equation}
\frac{2}{\sqrt{5\cdot4096}}
\approx
0.014.
\end{equation}
If the final MW estimate averages over \(N=20\) independent random parameter
samples, the shot-noise contribution to the mean is reduced by \(\sqrt{N}\):
\begin{equation}
\frac{0.014}{\sqrt{20}}
\approx
0.003.
\end{equation}

This justifies \(4096\) shots as a conservative upper-budget setting: finite
shot noise becomes small compared with the Monte Carlo variation over random
ansatz parameters and with possible hardware drift.

\section{Pilot Rule for Choosing Shots}

The pilot evaluates a grid of shot counts using the same ansatzes, depths,
random seed, and pilot parameter samples:
\begin{equation}
S \in \{512,1024,2048,4096\}.
\end{equation}
For each ansatz \(a\), depth \(L\), and consecutive shot pair
\((S_{k-1},S_k)\), compute
\begin{equation}
\Delta_S(a,L)
=
\left|
\bar Q_{a,L,S_k}
-
\bar Q_{a,L,S_{k-1}}
\right|.
\end{equation}
The selected shot count is the smallest \(S_k\) satisfying
\begin{equation}
\max_{a,L}\Delta_S(a,L)
\leq
\epsilon_S.
\label{eq:shot_rule}
\end{equation}
The default tolerance is
\begin{equation}
\epsilon_S = 0.02
\end{equation}
MW units. If no tested shot count satisfies Eq.~\eqref{eq:shot_rule}, the
protocol uses the largest tested shot budget.

The runner writes:
\begin{itemize}
  \item \texttt{shot\_pilot\_summary.csv},
  \item \texttt{shot\_stability.csv},
  \item \texttt{shot\_stability\_aggregate.csv}.
\end{itemize}

\section{Sample-Count Justification}

Once shots are fixed, the dominant statistical uncertainty is Monte Carlo
variation over random ansatz parameter vectors. For each ansatz/depth
configuration, let \(Q_1,\dots,Q_N\) be MW scores from \(N\) random parameter
samples. The sample mean is
\begin{equation}
\bar Q
=
\frac{1}{N}
\sum_{j=1}^{N}Q_j,
\end{equation}
with empirical standard deviation \(s_Q\). The estimated 95\% confidence
half-width is
\begin{equation}
h_N
=
z_{0.975}
\frac{s_Q}{\sqrt{N}},
\label{eq:half_width}
\end{equation}
where the implementation uses \(z_{0.975}=1.96\) by default.

The pilot chooses the smallest sample count \(N\) satisfying
\begin{equation}
\max_{a,L} h_N(a,L)
\leq
h_{\mathrm{target}},
\label{eq:sample_rule}
\end{equation}
with default target
\begin{equation}
h_{\mathrm{target}}=0.03
\end{equation}
MW units.

Equivalently, for a pilot-estimated standard deviation \(s_Q\), the required
number of samples is
\begin{equation}
N_{\mathrm{req}}
=
\left\lceil
\left(
\frac{1.96\,s_Q}{h_{\mathrm{target}}}
\right)^2
\right\rceil.
\label{eq:n_req}
\end{equation}
For example, if \(s_Q\leq0.06\) and \(h_{\mathrm{target}}=0.03\), then
\begin{equation}
1.96\frac{0.06}{\sqrt{20}}
\approx
0.026,
\end{equation}
so \(N=20\) meets the target half-width.

The runner writes:
\begin{itemize}
  \item \texttt{sample\_pilot\_summary.csv},
  \item \texttt{sample\_precision.csv},
  \item \texttt{sample\_precision\_aggregate.csv}.
\end{itemize}

\section{Hardware Iterations and Drift}

The shot and sample pilots quantify statistical precision, but not
time-dependent hardware drift. To estimate drift, repeat the frozen protocol with
the same random seed and choose the smallest iteration count \(K\) whose
worst-case drift half-width meets a target.

For repeated hardware sweeps \(\bar Q_1,\bar Q_2,\dots,\bar Q_K\) at fixed shots
and \(N\), define the iteration standard deviation \(s_{\mathrm{iter}}\) and
Student-\(t\) 95\% half-width
\begin{equation}
h_{\mathrm{iter}}
=
t_{0.975,K-1}\,
\frac{s_{\mathrm{iter}}}{\sqrt{K}}.
\end{equation}
Choose the smallest \(K \ge K_{\min}\) such that
\begin{equation}
\max_{\text{ansatz, depth}} h_{\mathrm{iter}} \le h_{\mathrm{target}}.
\end{equation}
Defaults are \texttt{--min-iterations 3}, \texttt{--max-iterations 5}, and
\texttt{--target-iteration-half-width 0.01}. Use \texttt{0.02} for qualitative
depth trends.

After shots and samples are already frozen, run only the drift pilot:
\begin{verbatim}
python scripts/run_iqm_mw_pilot.py \
  --drift-only \
  --shots 1024 \
  --samples 60 \
  --target-iteration-half-width 0.01 \
  --min-iterations 3 \
  --max-iterations 5 \
  --pilot-id mw_drift_pilot_final
\end{verbatim}

Using the same seed fixes the random parameter vectors across iterations. Thus,
iteration-to-iteration differences primarily reflect QPU drift, calibration
changes, queue conditions, and residual shot noise rather than different random
points in parameter space.

The runner writes:
\begin{itemize}
  \item \texttt{iteration\_summary.csv},
  \item \texttt{iteration\_stability.csv},
  \item \texttt{iteration\_precision.csv},
  \item \texttt{iteration\_precision\_aggregate.csv}.
\end{itemize}

\section{Recommended Paper-Quality Pilot}

A practical paper-quality pilot is:
\begin{verbatim}
python scripts/run_iqm_mw_pilot.py \
  --pilot-id mw_full_precision_pilot \
  --shot-grid 512 1024 2048 4096 \
  --sample-grid 10 20 40 60 \
  --pilot-samples 10 \
  --shot-tolerance 0.02 \
  --target-half-width 0.03 \
  --target-iteration-half-width 0.01 \
  --min-iterations 3 \
  --max-iterations 5
\end{verbatim}

For the current two-ansatz, three-depth setup, one full MW sweep costs
\begin{equation}
2 \times 3 \times N \times 3
\end{equation}
circuits, because each random parameter sample requires \(Z\), \(X\), and
\(Y\)-basis measurements. At \(N=20\), this is
\begin{equation}
2 \times 3 \times 20 \times 3 = 360
\end{equation}
circuits per full final sweep.

The recommendation file
\begin{verbatim}
mw_protocol_recommendation.json
\end{verbatim}
records the chosen shots, chosen sample count, chosen iterations,
target iteration half-width, min/max iteration bounds, iteration target flag,
shot-noise bound, mean-level shot-noise bound, and the exact decision rules
used by the pilot.

\section{References}

\begin{itemize}
  \item D. A. Meyer and N. R. Wallach, ``Global entanglement in multiparticle
  systems,'' \emph{Journal of Mathematical Physics}, 2002.
  \item M. A. Nielsen and I. L. Chuang, \emph{Quantum Computation and Quantum
  Information}, for Bloch-vector representation and Pauli measurements.
  \item M. Paris and J. Rehacek, \emph{Quantum State Estimation}, 2004, for
  quantum tomography methodology.
  \item W. Hoeffding, ``Probability inequalities for sums of bounded random
  variables,'' 1963, for concentration of bounded measurement outcomes.
\end{itemize}

\end{document}
